% Part: many-valued-logic
% Chapter: infinite-valued-logics
% Section: lukasiewicz

\documentclass[../../../include/open-logic-section]{subfiles}

\begin{document}

\olfileid{mvl}{inf}{luk}

\olsection{منطق~\L ukasiewicz}

\begin{editorial}
  ليس هذا إلا «بذرة» وجيزة لقسم في منطق~\L ukasiewicz اللانهائي القيم.
\end{editorial}

\begin{defn}\ollabel{def:lukasiewicz} مصفوفة منطق~\L ukasiewicz
اللانهائي القيم~$\LogLuk[\infty]$ هي:
\begin{enumerate}
  \item الجزء الخالي من الثوابت~$\Lang L_\OLMathNumeral{0}^{-}$ من لغة القضايا القياسية، ذو الروابط
  $\lnot$ و$\land$ و$\lor$ و$\lif$.
  % تصحيح مصدر مشترك (MVL-L0-I1): حُددت اللغة الخالية من lfalse التي تفسر المصفوفة جميع روابطها.
  \item مجموعة قيم الصدق~$\OLId{latin-looped}{V}_\infty$.
  \item القيمة المميزة وحدها هي~$\OLMathNumeral{1}$؛ أي~$\OLId{latin-looped}{V}^+ = \{\OLMathNumeral{1}\}$.
  \item دوال الصدق هي الدوال الآتية:
  \begin{align*}
    \tf{\lnot}[\LogLuk](\OLId{latin-ordinary}{x}) & = \OLMathNumeral{1} - \OLId{latin-ordinary}{x}\\
    \tf{\land}[\LogLuk](\OLId{latin-ordinary}{x},\OLId{latin-ordinary}{y}) & = \min(\OLId{latin-ordinary}{x},\OLId{latin-ordinary}{y})\\
    \tf{\lor}[\LogLuk](\OLId{latin-ordinary}{x},\OLId{latin-ordinary}{y}) & = \max(\OLId{latin-ordinary}{x},\OLId{latin-ordinary}{y})\\
    \tf{\lif}[\LogLuk](\OLId{latin-ordinary}{x},\OLId{latin-ordinary}{y}) & = \min(\OLMathNumeral{1},\OLMathNumeral{1}-(\OLId{latin-ordinary}{x}-\OLId{latin-ordinary}{y})) = \begin{cases}
      \OLMathNumeral{1} & \text{إذا كان } \OLId{latin-ordinary}{x} \le \OLId{latin-ordinary}{y}\\
      \OLMathNumeral{1}-(\OLId{latin-ordinary}{x}-\OLId{latin-ordinary}{y}) & \text{في غير ذلك.}
    \end{cases}
    \end{align*}
\end{enumerate}
وأما منطق~\L ukasiewicz ذو~$\OLId{latin-ordinary}{m}$ من القيم فيعرف على الوجه نفسه، إلا
أن~$\OLId{latin-looped}{V} = \OLId{latin-looped}{V}_\OLId{latin-ordinary}{m}$.
\end{defn}

\begin{prop}
  المنطق~$\LogLuk[\OLMathNumeral{3}]$ المعرف في \olref[thr][luk]{def:lukasiewicz}
  هو بعينه~$\LogLuk[\OLMathNumeral{3}]$ المعرف في \olref{def:lukasiewicz}.
\end{prop}

\begin{proof}
  يتبين ذلك إذا قابلنا جداول الروابط في
  \olref[thr][luk]{def:lukasiewicz} بجداول الصدق المستخرجة من
  معادلات \olref{def:lukasiewicz}:
  \begin{center}
    \begin{tabular}{c|c}
      $\tf{\lnot}$ & \\
      \hline
      $\OLMathNumeral{1}$ & $\OLMathNumeral{0}$ \\
      $\OLMathNumeral{1}/\OLMathNumeral{2}$ & $\OLMathNumeral{1}/\OLMathNumeral{2}$ \\
      $\OLMathNumeral{0}$ & $\OLMathNumeral{1}$
    \end{tabular}
    \quad
    \begin{tabular}{c|ccc}
      $\tf{\land}[\LogLuk[\OLMathNumeral{3}]]$ & $\OLMathNumeral{1}$ & $\OLMathNumeral{1}/\OLMathNumeral{2}$ & $\OLMathNumeral{0}$ \\
      \hline
      $\OLMathNumeral{1}$ & $\OLMathNumeral{1}$ & $\OLMathNumeral{1}/\OLMathNumeral{2}$ & $\OLMathNumeral{0}$ \\
      $\OLMathNumeral{1}/\OLMathNumeral{2}$ & $\OLMathNumeral{1}/\OLMathNumeral{2}$ & $\OLMathNumeral{1}/\OLMathNumeral{2}$ & $\OLMathNumeral{0}$\\
      $\OLMathNumeral{0}$ & $\OLMathNumeral{0}$ & $\OLMathNumeral{0}$ & $\OLMathNumeral{0}$
    \end{tabular}
    \\[2ex]
    \begin{tabular}{c|ccc}
      $\tf{\lor}[\LogLuk[\OLMathNumeral{3}]]$ & $\OLMathNumeral{1}$ & $\OLMathNumeral{1}/\OLMathNumeral{2}$ & $\OLMathNumeral{0}$ \\
      \hline
      $\OLMathNumeral{1}$ & $\OLMathNumeral{1}$ & $\OLMathNumeral{1}$ & $\OLMathNumeral{1}$ \\
      $\OLMathNumeral{1}/\OLMathNumeral{2}$ & $\OLMathNumeral{1}$ & $\OLMathNumeral{1}/\OLMathNumeral{2}$ & $\OLMathNumeral{1}/\OLMathNumeral{2}$ \\
      $\OLMathNumeral{0}$ & $\OLMathNumeral{1}$ & $\OLMathNumeral{1}/\OLMathNumeral{2}$ & $\OLMathNumeral{0}$
    \end{tabular}
    \quad
    \begin{tabular}{c|ccc}
      $\tf{\lif}[\LogLuk[\OLMathNumeral{3}]]$ & $\OLMathNumeral{1}$ & $\OLMathNumeral{1}/\OLMathNumeral{2}$ & $\OLMathNumeral{0}$ \\
      \hline
      $\OLMathNumeral{1}$ & $\OLMathNumeral{1}$ & $\OLMathNumeral{1}/\OLMathNumeral{2}$ & $\OLMathNumeral{0}$ \\
      $\OLMathNumeral{1}/\OLMathNumeral{2}$ & $\OLMathNumeral{1}$ & $\OLMathNumeral{1}$ & $\OLMathNumeral{1}/\OLMathNumeral{2}$  \\
      $\OLMathNumeral{0}$ & $\OLMathNumeral{1}$ & $\OLMathNumeral{1}$ & $\OLMathNumeral{1}$
    \end{tabular}
  \end{center}
\end{proof}

\begin{prop}\ollabel{prop:luk-infty-m}
  إذا كان~$\OLId{greek-variable}{Gamma} \Entails[\LogLuk[\infty]] !B$، كان
  $\OLId{greek-variable}{Gamma} \Entails[\LogLuk[\OLId{latin-ordinary}{m}]] !B$ لكل~$\OLId{latin-ordinary}{m} \ge \OLMathNumeral{2}$.
\end{prop}

\begin{proof}
  إثبات ذلك تمرين.
\end{proof}

\begin{prob}
  أثبت \olref[mvl][inf][luk]{prop:luk-infty-m}.
\end{prob}

% تصحيح مصدر (0400-I1): لا يصح إطلاق العكس لمجموعة مقدمات غير منتهية.
وأما عكس هذا الحكم فيصح إذا كانت مجموعة المقدمات متناهية.

ومنطق~\L ukasiewicz اللانهائي القيم أشهر المنطقات الضبابية. وأما
الشرط في أدبيات المنطق الضبابي فكثيرًا ما يعرف بـ~$\lnot !A \lor !B$؛
والحاصل حينئذ منطق كليني قوي لانهائي القيم.

\begin{prob}
  بيّن أن~$(\OLId{latin-ordinary}{p} \lif \OLId{latin-ordinary}{q}) \lor (\OLId{latin-ordinary}{q} \lif \OLId{latin-ordinary}{p})$ تحصيل منطقي في
  $\LogLuk[\infty]$.
\end{prob}

\end{document}
